\section{Affective-Precision Update and Policy Selection}
\label{app:affective_update}

This appendix specifies the auxiliary affective-precision update layered on top of the focal agent's partner-local POMDPs. The update maps partner-response model fit into a categorical posterior over inverse precision \(\beta_k\), then maps the posterior mean to partner-specific policy precision \(\gamma_k\). The same appendix also specifies how \(\gamma_k\) enters partner-choice policy comparison.

\subsection{Categorical beta support}

Each partner maintains a categorical posterior $q(\beta_k)$ over discrete inverse-precision levels $\{0.5, 0.67, 1.0, 1.5, 2.0\}$, updated each round from affective charge $\phi_k$. The categorical representation keeps the inverse-precision convention explicit, bounds the precision dynamics, and allows the update to be analyzed as movement over interpretable confidence states.

\subsection{Persistence prior}

The update proceeds in three steps. First, a persistence prior smooths the previous posterior via a tridiagonal transition matrix:
\begin{equation}
p(\beta_k) \leftarrow T \cdot q_{\text{prev}}(\beta_k),
\label{eq:persistence}
\end{equation}
where $T$ is tridiagonal with persistence parameter 0.8 and reflecting boundaries, encoding the prior expectation that precision evolves slowly.

\subsection{Charge-based pseudo-likelihood}

Second, a charge-based pseudo-likelihood maps positive charge toward high policy precision (low $\beta$) and negative charge toward low policy precision (high $\beta$):
\begin{equation}
\log \ell(\beta_\ell) = \phi_k \cdot \frac{1}{\beta_\ell},
\label{eq:pseudo_lik}
\end{equation}
where $\beta_\ell$ ranges over the five support points. This auxiliary update satisfies the intended monotonicity: positive charge reinforces lower $\beta$, negative charge reinforces higher $\beta$, and the effect scales with effective precision at each support point.

\subsection{Posterior update and policy precision}

Third, the posterior is normalized:
\begin{equation}
q(\beta_k) \propto \exp(\log \ell(\beta_\ell)) \cdot p(\beta_k).
\label{eq:posterior_update}
\end{equation}
Policy precision is then set from the posterior mean as given in Eq.~\eqref{eq:gamma_update} of the main text. The quantity $\beta_k$ is inverse precision, so higher posterior mean $\beta_k$ lowers policy precision $\gamma_k$, while lower posterior mean $\beta_k$ raises $\gamma_k$. As noted in Appendix~\ref{app:pomdp_beta_auxiliary}, $\beta_k$ is maintained as an auxiliary tracker outside the POMDP.

\subsection{Cross-partner policy selection}

In partner-choice regimes, partner-specific precision sharpens or flattens policy commitment among policies directed toward a given partner while preserving the mean policy evidence associated with that partner. The transformation in Eq.~\eqref{eq:centered_policy_scores} applies $\gamma_k$ to deviations from the partner-specific policy mean, rather than to raw policy scores. The agent samples or selects policies from the combined set of transformed scores across all partners and policies. This lets $\gamma_k$ control commitment among policies involving partner $k$ without changing the average evidence for approaching partner $k$ during cross-partner comparison.

\subsection{Ablation variants}

Runtime variants are summarized in Table~\ref{tab:precision_variants} (see also Section~\ref{sec:simulation_setup}).

\Needspace{10\baselineskip}
\begin{table}[H]
\centering
\small
\caption{Precision-runtime ablation variants.}
\label{tab:precision_variants}
\begin{tabularx}{\linewidth}{@{}>{\raggedright\arraybackslash}p{0.18\linewidth}X>{\raggedright\arraybackslash}p{0.24\linewidth}@{}}
\toprule
Variant & $\beta$ tracker and $\gamma_k$ & Precision scope \\
\midrule
No-affect & off; $\gamma_k = \gamma_{\text{base}}$ fixed & --- \\
Tracked-only & $q(\beta_k)$ updated; $\gamma_k = \gamma_{\text{base}}$ fixed & per-partner \\
Partner-local & $q(\beta_k)$ updated; $\gamma_k$ from $q(\beta_k)$ & per-partner \\
Shared-$\beta$ & pooled $q(\beta)$ updated; $\gamma_k$ from shared $\beta$ & pooled tracker; separate beliefs \\
\bottomrule
\end{tabularx}
\end{table}
